% Privacy-first section 8: limitations and discussion. Budget 0.75 page, no figure.
\section{Limitations and discussion}
\label{sec:discussion}

\paragraph{Cohort size is held fixed, and we do not know its effect.}
Every result here uses one fine-tuning corpus size. Scaling the number of individuals, or the amount
of sequence per individual, plausibly changes all of it: a planted record is a smaller fraction of
the gradient signal in a larger corpus, which argues for less memorisation, while a larger cohort
contains more rare haplotypes to confuse an identification probe, which argues for less
identification but also for a harder task. We tested neither direction. This is the first
experiment we would run next; it is cheap on the controlled backbone and expensive on released
weights.

\paragraph{The tokenisation result is scoped, and one released model sits outside the scope.}
Table~\ref{tab:slots} varies $k$ with the architecture, objective, optimiser and corpus fixed, which
is the only setting in which the axis is clean. It requires a model we train ourselves, because no
released family offers it. Across architectures the ordering does not survive: HyenaDNA, a released
single-nucleotide model at 14.2M parameters, almost exactly the size of our controlled backbone,
reaches 95\% verbatim extraction \citep{nguyen2023hyenadna}. That comparison varies architecture,
pretraining objective and objective-match simultaneously, so it does not contradict
\S\ref{sec:slots}; it bounds its generality. ``Finer tokenisation is safer'' is false as a general
claim and we do not make it.

\paragraph{The human arm is one chromosome.}
Canaries are drawn from chromosome~22 of a single cohort. Chromosome~21 appears only as pretraining
and replication. Nothing in the construction is chromosome-specific, but we have not shown that.

\paragraph{Our decoders are lower bounds, including where they read zero.}
We report whichever decoder is stronger per model, so every extraction figure understates what a
better attack could obtain. Two of our zeros are demonstrably decoder limits rather than model
properties: on DNABERT-2 the leave-one-out ceiling for single-copy records is 0.21--0.39 and on the
overlapping DNABERT family it is 0.70--0.87, while our searches return nothing in both cases. A reader
should treat those zeros as ``not reached by this attack'' and nothing more. From the same kind of
zero, the field currently draws the opposite inference.

\paragraph{Extraction rates carry more uncertainty than the probe count suggests.}
The re-masking decode is a discrete iteration and is sensitive to the model's last-bit arithmetic.
Re-running an identical cell after the host's floating-point behaviour changed moved its extraction
rate by six points, comparable to the binomial standard error on 50 probes (\S\ref{sec:measure}).
Identification does not share this sensitivity. Anyone reproducing these numbers should expect
agreement on identification to the digit and on extraction to within several points.

\paragraph{What we think this changes.}
The practical recommendation, which is to stop at the held-out optimum, is sound and nearly free,
but it rests on a quantity nobody in this field tracks, and the margin is three hundredths of a bit.
The second point concerns what a null result licenses. A team that plants canaries, measures extraction, reads
zero at single copy and ships has done what the literature currently asks, and has shipped a model
that will tell an adversary which of their participants it was trained on. Releasing fine-tuned
genomic weights should require reporting both numbers, measured with a decoder matched to the model's
own tokenizer, with a never-trained control scored in the same run. None of that is expensive, and
all of it is currently optional.
